\documentclass[11pt,a4paper]{article}
\usepackage[utf8]{inputenc}
\usepackage[T1]{fontenc}
\usepackage{amsmath, amssymb, amsfonts}
\usepackage{graphicx}
\usepackage{geometry}
\geometry{a4paper, margin=1in}
\usepackage{hyperref}
\usepackage{xcolor}
\usepackage{listings}
\usepackage{tikz}
\usetikzlibrary{shapes.geometric, arrows, positioning}

\title{\textbf{Understanding smolVLA: Architecture and Data Pipeline}}
\author{Explaining the Vision-Language-Action Model}
\date{\vspace{-5ex}}

\begin{document}

\maketitle

\section{Introduction to smolVLA}
The \textbf{smolVLA} model is a compact, efficient Vision-Language-Action (VLA) architecture designed for affordable robotics. Unlike traditional massive VLAs with billions of parameters, smolVLA is optimized to be trained on a single GPU and deployed on consumer-grade hardware or CPUs. It achieves this efficiency by adapting a small pretrained Vision-Language Model (VLM) and coupling it with a lightweight Action Expert that generates continuous control actions.

At a high level, the model receives three primary sensory modalities from the robot's environment, processes them to form a cohesive contextual representation, and then generates a sequence (or ``chunk'') of low-level actions.

\section{The Architecture: VLM + Action Expert}
SmolVLA fundamentally consists of two components interconnected via attention mechanisms:

\begin{enumerate}
    \item \textbf{The Pretrained Vision-Language Model (VLM):} Based on the \texttt{SmolVLM-2} architecture (using \texttt{SigLIP} for vision and a 30-layer language decoder). To save computational resources, smolVLA \textit{truncates the LLM backbone}. Instead of passing through all 30 layers, the model is configured to use only the first $N$ layers (typically $N=16$). These intermediate layers are found to contain sufficiently rich semantic information for robotic control, while effectively halving the computation cost of the VLM forward pass.
    
    \item \textbf{The Action Expert:} A conditional Flow Matching Transformer that predicts continuous action chunks. Because smolVLA is an extension of a decoder-only architecture (Llama-style), it does not use separate cross-attention layers. Instead, the Action Expert simply relies on the standard \textbf{Self-Attention layers} applied to the concatenated sequence \texttt{[Prefix, Suffix]}. This allows the action suffix tokens to organically ``attend to'' the rich contextual representations generated by the VLM prefix, predicting the exact velocity vector needed to denoise the chunk.
\end{enumerate}

\subsection{The Pretrained Backbones}
To achieve its high efficiency, smolVLA stands on the shoulders of two specialized "foundational" models:

\begin{itemize}
    \item \textbf{Vision: SigLIP (Vision Transformer):} For visual perception, smolVLA uses a **SigLIP encoder**. Unlike older vision models, SigLIP uses a unique \textbf{Sigmoid Loss} during training, which allows it to learn more precise connections between images and text. Structurally, it is an \textbf{Encoder-only ViT} that treats the image as a grid of $32 \times 32$ patches, extracting fine-grained features from every corner of the robot's workspace.
    \item \textbf{Language: SmolLM2 (Causal Decoder):} Instruction understanding is handled by the first 16 layers of **SmolLM2-500M**. This is a **Llama-style causal transformer** optimized for tiny footprints. Its primary job is to turn raw sub-word tokens into a "contextual history" that the Action Expert can cross-reference to understand the specific intent of the human's command.
\end{itemize}

\section{The Input Pipeline: Preparing the Prefix}
Before reaching the transformer architecture, the three modalities (Vision, Language Instruction, and Robot State) are heavily preprocessed, vectorized, and concatenated.

\subsection{Vision Data Processing}
\textbf{Python Implementation:} \texttt{processor\_smolvla.py} (Preprocessing), \texttt{smolvlm\_with\_expert.py} (Encoder/Projection)

The robot's visual observations (incoming from one or more RGB cameras) form a temporal video stream. However, to maintain efficiency, smolVLA processes this data with specific simplifications:
\begin{itemize}
    \item \textbf{Temporal Dimension (Time):} Although the raw dataset provides a continuous video (a 5D tensor of shape $\text{Batch} \times \text{Time} \times \text{Channel} \times \text{Height} \times \text{Width}$), smolVLA \textbf{discards the temporal history} of the images. It extracts only the most recent frame (the last index in the time dimension) to use as the current visual observation. This means the model acts based purely on the current visual snapshot.
    \item \textbf{Resizing and Padding:} The extracted current frame is resized to match the strict input resolution expected by the vision encoder ($512 \times 512$). Because the raw camera aspect ratios usually do not match this square resolution, the image is scaled proportionally, and symmetric padding (with a value of $0$) is added to fill the remaining space without distorting the objects in the image.
    \item \textbf{Normalization:} The raw pixel values usually arrive in the range $[0, 1]$. To match the distribution expected by the pretrained \texttt{SigLIP} vision encoder, the image tensor is arithmetically transformed by multiplying by 2 and subtracting 1. This mathematical operation is broadcasted across the tensor, applying equally to the Red, Green, and Blue channels, shifting all pixel values to the $[-1, 1]$ range.
    \item \textbf{Tokenization (Pixel Shuffle):} The normalized images are processed by the encoder into a grid of $1,024$ patches ($32 \times 32$). To reduce computational load, smolVLA applies a \textbf{Pixel Shuffle (Space-to-Depth)} operation. Unlike a simple average which would "blur" information, this operation takes a $4 \times 4$ local neighborhood of patches (16 total) and \textbf{concatenates (stacks)} their feature vectors end-to-end. This is a deliberate trade-off: we reduce the sequence length from 1,024 to just \textbf{64 visual tokens}, but each token becomes 16 times "deeper" in information.
    \begin{itemize}
        \item \textbf{Overcoming the Quadratic Wall ($L^2$):} The computational cost of Self-Attention is \textbf{quadratic} in terms of the number of tokens ($L$). By shuffling 16 short vectors into 1 long vector, the sequence length is reduced by 16x. Because of the $L^2$ scaling, this reduces the total attention operations by \textbf{256x} ($16^2$), making the model significantly faster and more memory-efficient without losing raw data.
        \item \textbf{The "Goldilocks" Trade-off (Why not 1 Token?):} While fewer tokens are faster, compressing the entire image into just \textbf{1 token} would create a massive \textbf{Information Bottleneck} (trying to fit all details into just 960 dimensions). Furthermore, the \textbf{Attention Mechanism} requires multiple tokens to compare relationships (e.g., comparing the robot's gripper to the object). The \textbf{64 tokens} maintain an $8 \times 8$ grid, preserving enough \textbf{Spatial Awareness} for the model to know *where* objects are relative to each other, while still bypassing the quadratic math explosion.
    \end{itemize}
    \item \textbf{Frozen Model Projection:} After the stacking, these deep vectors are passed through a \textbf{frozen linear projection} (resampler) that compresses them back to the model's standard hidden dimension of \textbf{960}. Because this projector is part of the original pretrained VLM, it has already learned a sophisticated "summary" function to preserve the most relevant visual details (like object edges and textures) within these 960 dimensions.
    \item \textbf{Multi-Camera Concatenation:} If the model uses multiple cameras (e.g., a "center" and a "wrist" camera), each image is processed independently through the same pipeline. The resulting 64-token sequences are simply \textbf{concatenated end-to-end}. For a two-camera setup, this results in a visual prefix of 128 tokens total before appending language and state data.
    \item \textbf{Handling Missing Cameras:} If the policy expects up to 3 cameras (by default, it expects 3 cameras) but only 1 is active, the model generates fully padded ``dummy'' images filled with $-1$ to maintain a consistent tensor shape.
\end{itemize}

\textbf{Result of the Vision Pipeline:} For each camera observation at time $t$, the vision pipeline transforms a $512 \times 512$ RGB image into a sequence of \textbf{64 visual tokens}. Each token is a high-dimensional vector of length \textbf{960}. In a multi-camera setup (e.g., center and wrist), these sequences are simply appended together, forming the first part of the model's contextual memory. 

\subsection{Language Instruction Processing}
\textbf{Python Implementation:} \texttt{processor\_smolvla.py} (Tokenization), \texttt{modeling\_smolvla.py} (Embeddings)

The language instruction (e.g., \textit{``Grasp the object''}) is processed in three distinct phases:
\begin{itemize}
    \item \textbf{Phase 1: Tokenization (Discrete IDs):} The raw string is converted into a list of \textbf{discrete integer IDs} using the \texttt{SmolVLM2} dictionary. Note that the model uses \textit{sub-word tokenization}: common words like "the" are 1 token, but complex words like "grasping" may be split into two pieces ("grasp" + "ing"). Additionally, a mandatory **newline token** (\verb|\n|) is always appended to the end. Consequently, a 5-word sentence often results in 7 or 8 tokens.
    \item \textbf{Phase 2: Static Embedding (The Lookup Table):} These integer IDs are used to index into a giant \textbf{Embedding Matrix}. For each ID, the model retrieves a "static" 960-dimensional vector.
    \item \textbf{Phase 3: Contextualization (The Transformer Pass):} These static vectors are concatenated with the visual tokens and passed together through the $N$ layers of the \textbf{frozen transformer backbone}. 
    \begin{itemize}
        \item \textbf{Multimodal Interaction:} Because they share the same attention space, language tokens can "attend" to visual tokens (and vice versa) during this pass. For example, the vector for the word "grasp" can mathematically align itself with the specific visual patches containing the robot's gripper or the target object \textit{before} the information ever reaches the Action Expert.
    \end{itemize}
\end{itemize}

\textbf{Result of the Language Pipeline:} The instruction starts as text and ends as a sequence of \textbf{contextualized 960-dimensional vectors}. The number of vectors corresponds to the number of tokens (usually $\text{words} + 2 \text{ or } 3$), each having a length of \textbf{960}.

\subsection{Robot State (Proprioception) Processing}
\textbf{Python Implementation:} \texttt{processor\_smolvla.py} (Stats), \texttt{modeling\_smolvla.py} (Projector)

To make informed decisions, smolVLA must know the current physical configuration of the robot (its "proprioception"). Unlike vision and language, which use complex frozen backbones, the robot state is integrated via a lightweight, learned projection:
\begin{itemize}
    \item \textbf{Input Padding:} Different robots have different numbers of sensors (e.g., a 6-DOF arm vs. a 14-DOF bimanual robot). To maintain a consistent architecture, smolVLA zero-pads every state vector up to a fixed maximum dimension (typically \texttt{max\_state\_dim = 32}). If your robot has only 4 joints, the remaining 28 slots are filled with $0$s.
    \item \textbf{The State Projector (\texttt{state\_proj}):} This padded vector is passed through a single \textbf{Trainable Linear Layer} (\texttt{nn.Linear}). Unlike the frozen VLM backbone, this specific layer is actively trained. 
    \begin{itemize}
        \item \textbf{Pretrained Initialization:} When loading the base model, this layer comes initialized with pretrained weights. It already "knows" how to map physical states into the model's 960-dimensional space.
        \item \textbf{Fine-Tuning adaptation:} Because the configuration sets \texttt{train\_state\_proj = True}, these weights continue to update during fine-tuning. This allows the model to leverage its pretrained robotics knowledge while slightly adjusting to perfectly match the mechanics of a new specific robot arm or task.
    \end{itemize}
    \item \textbf{Linear Transformation:} Interestingly, this is a pure linear transformation (\(y = xW + b\)) \textit{without} a non-linear activation function (like ReLU or SiLU). A single linear layer is sufficient to project the continuous physical values into the high-dimensional hidden space.
\end{itemize}

\textbf{Result of the State Pipeline:} The physical state of the robot is compressed into a \textbf{single state token} (one vector) of length \textbf{960}. This single vector follows the visual and language tokens, completing the prefix sequence. 

\subsection{Summarizing the Prefix Sequence}
At the end of these three independent pipelines, smolVLA concatenates the resulting vectors along the sequence dimension to form a unified \textbf{Prefix}. This prefix serves as the comprehensive "memory" or "context" that conditions the robot's future actions.

\textbf{Composition Details:}
\begin{itemize}
    \item \textbf{Concatenation Order:} The vectors are simply appended end-to-end: $[ \text{Vision} , \text{Language} , \text{State} ]$.
    \item \textbf{The Resulting Sequence:} If an instruction of $n$ words is tokenized into $n+2$ tokens (accounting for sub-words and the mandatory newline), the sequence looks like:
    \[ (v_1, \dots, v_{64}, \underbrace{l_1, \dots, l_{n+2}}_{\text{Instruction Tokens}}, s_{proprio}) \]
    \item \textbf{Uniform Dimension:} Every single vector in this sequence (whether it came from a camera patch, a word, or a joint sensor) is exactly \textbf{960} elements long.
\end{itemize}

\textbf{The Full Prefix sequence formula:}
\[ \text{Prefix} = [ \underbrace{\text{64 Visual Tokens}}_{\text{Frozen SigLIP + Shuffle}} , \underbrace{\text{N Language Tokens}}_{\text{Frozen SmolLM2}} , \underbrace{\text{1 State Token}}_{\text{Learned Projection}} ] \]

\textbf{Note on the "960" Dimension:}
While the codebase is flexible, the specific $SmolVLA-500M$ model uses a base width of \textbf{960} for all VLM-side tokens. This constant is not hard-coded but is dynamically loaded from the \texttt{hidden\_size} parameter of the pretrained configuration. All other dimensions are derived from this anchor: for example, the Action Expert's width of \textbf{720} is calculated as $960 \times 0.75$ (the \texttt{expert\_width\_multiplier}).

\section{The Action Suffix}
Once the prefix context is ready, the Action Expert calculates the physical motion. This is treated as a conditional flow-based generation task.

\subsection{Action Encoding (The Suffix)}
\textbf{Python Implementation:} \texttt{modeling\_smolvla.py} (\texttt{embed\_suffix})

The second part of the model's input is the \textbf{Suffix}, which represents a noisy estimate of the robot's future motion. Transforming raw numbers into the Expert's high-dimensional space requires two specific "upcasting" operations:

\begin{itemize}
    \item \textbf{Action Chunking (The 50 Steps):} smolVLA predicts an \textbf{Action Chunk} of typically \textbf{50 future timesteps} simultaneously.
    \begin{itemize}
        \item \textbf{Configuration Value:} This value is explicitly defined in the code as the default \texttt{chunk\_size = 50} within the \texttt{SmolVLAConfig} class. It dictates the length of the temporal window for which the model generates predicted actions in a single forward pass.
    \end{itemize}
    \item \textbf{Action Projection (\texttt{action\_in\_proj}):} Each raw action step is a 32-dimensional vector. To match the Action Expert's internal width (which is typically $75\%$ of the VLM's width), each step is passed through a \textbf{Learnable Linear Layer} that projects it from \textbf{32} dimensions up to \textbf{720} dimensions.
    \item \textbf{Absolute Values:} As a reminder, these are \textbf{absolute control values} (e.g., recorded as \textbf{87}$^\circ$, not a delta of $+2^\circ$).
    \item \textbf{Noise Interpolation (The Flow):} During training, we create a noisy sample $x_\tau$ by interpolating between the clean demonstration and random Gaussian noise using the parameter $\tau \in [0, 1]$.
    \[ x_\tau = \tau \cdot \text{Noise} + (1 - \tau) \cdot \text{Clean Action} \]
    \item \textbf{Sinusoidal Noise Level Embedding:} Simply feeding a scalar $\tau$ (like $0.2$) is insufficient for the network to distinguish fine noise levels. Instead, $\tau$ is transformed into a \textbf{Sinusoidal Positional Encoding}—a high-dimensional signature vector (length 720) composed of sine and cosine waves of varying frequencies.
    \item \textbf{Infinite Variety:} For a single physical observation (one prefix), we sample a random noise level $\tau$ for every single training batch. This means the model sees the same demonstration thousands of times, but each time with a different level of "messiness."
    \item \textbf{The Beta Distribution Trick:} In smolVLA, $\tau$ is sampled using a **Beta Distribution** ($\alpha=1.5, \beta=1.0$). This slightly biases the training toward higher noise levels ($\tau \approx 1$). This is intentional: the model needs to be extra "smart" at the very beginning of the flow when it only has pure noise to work with, as a mistake there would derail the entire trajectory.
    \item \textbf{The Suffix Fusion Pipeline:} For each of the 50 steps:
    \begin{enumerate}
        \item The 720-dim \textbf{Action Projection} and the 720-dim \textbf{Noise Embedding} are \textbf{concatenated} into a single 1440-dim vector.
        \item This large vector is passed through a \textbf{Trainable 2nd-stage MLP} (\texttt{action\_time\_mlp}).
        \item \textbf{Internal Non-linearity:} This MLP consists of two linear layers with a **SiLU (Swish)** activation function in between ($Linear \rightarrow SiLU \rightarrow Linear$).
        \item This MLP "squeezes" the combined information back into a final \textbf{720-dimensional Suffix Token}.
    \end{enumerate}
    \item \textbf{Action Normalization:} It is important to note that the raw sensor values (e.g., motor angles ranging from $-180^\circ$ to $180^\circ$) are \textit{not} fed directly into the noising process. Before any noise is added, the actions are \textbf{normalized} using the Mean and Standard Deviation statistics of the training dataset. This scales the typical range to approximately $[-1, 1]$, ensuring that the Gaussian noise ($\mathcal{N}(0, 1)$) and the actions exist on the same mathematical scale.
    \textbf{Note on SiLU (Swish):} The Sigmoid Linear Unit (\(\text{SiLU}(x) = x \cdot \sigma(x) = \frac{x}{1 + e^{-x}}\)) is used here instead of standard ReLU. Unlike ReLU, it is smooth, differentiable everywhere, and allows a small amount of negative information to flow, which leads to better convergence in large Transformer models.
\end{itemize}

\subsection{Prefix Conditioning and the KV Cache Optimization}
The Action Expert is essentially a continuation of the Transformer architecture. How does it actually ``condition'' the generation of the velocity vector on the images and text? 

\begin{itemize}
    \item \textbf{Sequence Concatenation:} The transformer is simply fed the sequence \texttt{[Prefix, Suffix]} glued together end-to-end. Because it is a Transformer, when it processes the \texttt{Suffix} (the noised actions), the self-attention heads allow the action tokens to organically ``attend to'' or ``look at'' the \texttt{Prefix} (the vision and text tokens). This means that when the model decides which direction to push the noise (the velocity vector), it intrinsically uses the contextual information it gathers directly from the images and text.
    \item \textbf{The KV Cache Optimization:} Flow Matching is an iterative process. During inference, the noised actions must be passed through the model multiple times to gradually denoise them. It would be prohibitively computationally expensive to re-process the heavy images and text every single time. Instead, the \textbf{Prefix (Images + Text + State)} is passed through the Vision-Language Model (VLM) \textit{only once}. The internal representations of the prefix at every transformer layer are saved in memory---this is called the \textbf{KV Cache} (Key-Value Cache). 
    \item \textbf{Efficient Iteration:} For every timestep of the denoising loop, only the \texttt{Suffix} (noised actions + time) is fed into the Action Expert part of the transformer. By attending to the saved KV Cache, the model instantly knows what the robot is looking at and what the instruction is, allowing it to rapidly predict the velocity vector iteratively.
\end{itemize}

\subsection{Final Adjustments: The "Pre-Layer" Pipeline}
\textbf{Python Implementation:} \texttt{modeling\_smolvla.py} (Scaling), \texttt{smolvlm\_with\_expert.py} (RoPE/Norm)

Just before the unified $X = (\text{Prefix}, \text{Suffix})$ enters the first layer of the Transformer, three final technical adjustments are applied to ensure numerical stability and spatial awareness:

\begin{itemize}
    \item \textbf{Embedding Scaling ($\sqrt{D}$):} The visual and language embeddings are multiplied by the square root of their dimension ($\sqrt{960} \approx 31$). This is a standard Transformer initialization trick that prevents the variance of the embeddings from collapsing as they are summed with positional information.
    \item \textbf{Position IDs (The "Sidecar" Argument):} Even though we have a sequence of vectors, the model needs to know their order (e.g., that $v_{10}$ is next to $v_{11}$). 
    \begin{itemize}
        \item \textbf{Simple Integers:} From the perspective of the input data, we simply prepare a list of integers: \texttt{position\_ids = [0, 1, 2, ..., 115]}. This is fed into the model as a separate argument (like a metadata "sidecar").
        \item \textbf{Instruction for the Model:} These integers are not added or concatenated to the vectors yet. Instead, they act as "pointers" that tell the internal Transformer layers exactly how to apply **Rotary Position Embeddings (RoPE)**. 
        \item \textbf{Internal Rotation:} Once inside the layers, the model uses these ID arguments to mathematically "twist" or rotate the vectors, allowing it to perceive both absolute and relative distance between items (like a Joint Angle and a task word) with perfect precision.
    \end{itemize}
    \item \textbf{Batching:} Everything is processed in parallel. If the batch size is 8, the model effectively sees 8 different robots and 8 different noisy trajectories simultaneously.
    \begin{itemize}
        \item \textbf{Massive Parallelism:} GPUs contain thousands of tiny computing units (CUDA cores). Instead of using one "gigantic calculator" for one robot, we assign groups of these cores to solve different robots in the batch at the exact same time. This completes 8 calculations in the same time as 1.
        \item \textbf{Stable Learning Direction:} During training, the model calculates the "direction" it needs to change its weights (the gradient). Gradients from a single demonstration can be erratic or noisy. By **averaging the loss** across all 8 samples in the batch, the model finds a smoother, more reliable direction that represents the general patterns of the whole dataset, preventing it from overreacting to outliers.
    \end{itemize}
    \item \textbf{First Layer Entrance (RMSNorm):} The very first operation inside Layer 0 of the Transformer is a **Root Mean Square Normalization (RMSNorm)**.
    \begin{itemize}
        \item \textbf{Rescaling, Not Centering:} Unlike standard LayerNorm, **RMSNorm does not subtract the mean.** It does not "center" the data at zero. Instead, it looks at each of the 115 vectors of size 960 independently and regulates their "energy" or "volume."
        \item \textbf{The Logic:} It calculates the Root Mean Square of the 960 elements and divides the vector by that value. This ensures the output vector has a consistent scale, preventing numerical values from "exploding" as they pass through deep layers.
        \item \textbf{Learned Gains:} After this rescaling, the model applies a set of \textbf{960 trainable parameters (gains)}, allowing it to learn exactly which feature dimensions (e.g., specific visual patterns or joint angles) should be boosted or suppressed for optimal decision making.
    \end{itemize}
\end{itemize}

\section{Action Normalization and Safety}
Because robot hardware uses specific physical units (like degrees for joints or meters for positions), and neural networks prefer small, stable numbers, smolVLA uses a strict normalization pipeline:

\subsection{The Normalization Mapping}
The model configuration (\texttt{SmolVLAConfig}) defines a \texttt{normalization\_mapping} where:
\begin{itemize}
    \item \textbf{Images:} Stay in their visual range (processed to $[-1, 1]$ for SigLIP).
    \item \textbf{State (Proprioception):} Normalized using \texttt{MEAN\_STD}.
    \item \textbf{Actions:} Normalized using \texttt{MEAN\_STD}.
\end{itemize}

\subsection{Why Normalize?}
If we were adding Gaussian noise (scale $\approx 1.0$) to a motor angle of $150^\circ$, the noise would be virtually invisible ($150 \pm 1$ is basically still $150$). By scaling the $150^\circ$ down to a representation like $0.8$ before adding noise, the noise becomes a significant signal that the model can learn to filter. This makes the Flow Matching mathematically well-behaved across different types of robots.

\section{The Training}

\subsection{Training Strategy: Learning the Field}
\textbf{Python Implementation:} \texttt{modeling\_smolvla.py} (\texttt{forward})

To make this flow work, the model must know the correct velocity for \textit{any} possible noise level. smolVLA achieves this by "augmenting" the data within the generative process itself.

\begin{itemize}
    \item \textbf{Batch \& Noise Matrix Generation:} Every single time the training loop grabs a batch of examples (e.g., $8$ examples), it instantaneously generates a \textbf{completely unique Noise Matrix} for each sample. Since an action chunk has $50$ timesteps and $32$ joints, this equates to $1,600$ independent random Gaussian values per robot observation.
    
    \item \textbf{Data Augmentation ($x_\tau$ and $\tau$):} For each instance in the batch, a unique scalar $\tau \in [0, 1]$ is drawn. The script then creates the corrupted chunk by blending the brand new Noise Matrix with the Clean Action:
    \[ \text{NoisedChunk}_{\tau} = \tau \cdot \text{NoiseMatrix} + (1 - \tau) \cdot \text{CleanChunk} \]
    This forces the model to learn the \textit{general mathematical rule} of de-noising rather than memorizing snapshots.

    \item \textbf{The Ground Truth Velocity ($u_\tau$):} Unlike traditional models where the "label" is a fixed category, the optimal velocity label is computed on the fly by subtracting the clean chunk from the newly generated Noise Matrix:
    \[ u_\tau = \text{NoiseMatrix} - \text{CleanChunk} \]
    \begin{itemize}
        \item \textbf{Grading the Model:} The Action Expert processes the Noised Chunk and $\tau$, and uses its cross-attention to the prefix to predict its own velocity $v_\tau$. We "grade" it using the \textbf{Mean Squared Error (MSE)} between its prediction $v_\tau$ and the true $u_\tau$. Once the weights are updated, that specific noise matrix and $u_\tau$ are discarded.
    \end{itemize}
\end{itemize}

\section{Model Inference}
At deployment time, the robot must generate a clean action trajectory starting from absolute chaos.

\subsection{1. The De-noising Flow}
\textbf{Python Implementation:} \texttt{modeling\_smolvla.py} (\texttt{inference})

At test time (inference), the robot doesn't know the answer, and it has no "ground truth" to interpolate with. Instead, it performs an iterative process called **Euler Integration**:
\begin{enumerate}
    \item \textbf{Initial State:} The robot starts with a suffix of pure random Gaussian noise ($\tau = 1$).
    \item \textbf{Querying the Expert:} It feeds the observations (\textbf{Prefix}) and the current noise (\textbf{Suffix}) to the model.
    \item \textbf{The Update:} The model predicts the **Velocity Vector** ($v_\tau$)—the direction that points toward a clean action.
    \item \textbf{Flowing:} The robot moves the noisy points a small step in that direction: $x_{new} = x_{old} + dt \cdot v_\tau$.
    \item \textbf{Repetition:} This is repeated for several small steps (usually **10 iterations**). By the time it reaches $\tau = 0$, the random noise has morphed into a smooth, executable action chunk.
\end{enumerate}

\section{Summary: The Data Flow}
\begin{enumerate}
    \item \verb|Observation Images| $\rightarrow$ Resize \& Pad $\rightarrow$ $[-1, 1]$ Normalize $\rightarrow$ SigLIP $\rightarrow$ \textbf{Vision Tokens}
    \item \verb|Task Instruction| $\rightarrow$ Append \verb|\n| $\rightarrow$ Tokenize $\rightarrow$ LLM Embedding $\rightarrow$ \textbf{Language Tokens}
    \item \verb|Robot Internal State| $\rightarrow$ Pad $\rightarrow$ Linear Projection $\rightarrow$ \textbf{State Tokens}
\end{enumerate}
\vspace{0.3cm}
$\Rightarrow$ Concatenate into \textbf{Prefix Sequence}. Forward pass through $N$ layers of VLM.
\vspace{0.3cm}
\begin{enumerate}
    \setcounter{enumi}{3}
    \item \verb|Noisy Actions| + \verb|Time Embedding| $\rightarrow$ MLP $\rightarrow$ \textbf{Suffix Sequence}
\end{enumerate}
$\Rightarrow$ Forward pass through Action Expert (Cross-Attending to Prefix). Optimize via MSE Flow Matching Loss.

\end{document}
